\documentclass[a4paper,twoside]{article}

\usepackage{morewrites}

\usepackage[svgnames,dvipsnames,x11names]{xcolor}
\usepackage{graphicx}
\usepackage[english]{babel}
\usepackage[colorlinks=true, linkcolor=NavyBlue, urlcolor=NavyBlue, citecolor=NavyBlue]{hyperref}
\usepackage[a4paper]{geometry}
\usepackage{ts-typesetting}
\usepackage{ts-fonts}
\usepackage{xcolor}
\usepackage{epigraph}
\usepackage{marginnote}
\usepackage{multicol}
\usepackage{progressbar}
\usepackage{graphicx}
\usepackage{tcolorbox}
\usepackage{fvextra}
\usepackage{amsmath}
\usepackage{float}
\usepackage{seqsplit}
\renewcommand*{\marginfont}{
\footnotesize
\sloppy
\emergencystretch=1em
\hyphenpenalty=50
\exhyphenpenalty=50
}
\newlength{\tsmarginparavail}
\newlength{\tsmarginparalt}
\newlength{\tsmarginparbuf}
\setlength{\tsmarginparbuf}{6mm}
\makeatletter
\AtBeginDocument{
\setlength{\tsmarginparavail}{\dimexpr\paperwidth-\textwidth-1in-\oddsidemargin-\marginparsep-\tsmarginparbuf\relax}
\ifdim\tsmarginparavail<\z@\setlength{\tsmarginparavail}{\z@}\fi
\if@twoside
\setlength{\tsmarginparalt}{\dimexpr1in+\evensidemargin-\marginparsep-\tsmarginparbuf\relax}
\ifdim\tsmarginparalt<\z@\setlength{\tsmarginparalt}{\z@}\fi
\ifdim\tsmarginparalt<\tsmarginparavail
\setlength{\tsmarginparavail}{\tsmarginparalt}
\fi
\fi
\ifdim\tsmarginparavail<\marginparwidth
\setlength{\marginparwidth}{\tsmarginparavail}
\fi
}
\makeatother
\usepackage{ts-code}

\usepackage{lastpage}
\usepackage{refcount}
\makeatletter
\def\ps@plain{
\let\@oddhead\@empty
\let\@evenhead\@empty
\def\@oddfoot{
\hfil
\ifnum\getpagerefnumber{LastPage}>1\relax
\thepage
\fi
\hfil}
\let\@evenfoot\@oddfoot
}
\makeatother

\title{Margin Notes\\[0.5em]\large The \tscodeinline{\{aside\}[…]} inline extension}
\author{TeXSmith}\date{}
\begin{document}
\maketitle
\section{Margin Notes}
TeXSmith's margin-note extension adds a single inline shorthand to the
unified \tscodeinline{\{keyword\}[content]} family already used by \tscodeinline{\{index\}[term]} and
\tscodeinline{\{raw latex\}(payload)}. The syntax is:
\begin{tscode}[lang=text, engine=pygments]
\begin{Verbatim}[commandchars=\\\{\},breaklines, breakanywhere, commandchars=\\\{\}]
\PYZob{}aside side=…\PYZcb{}[note text]
\end{Verbatim}
\end{tscode}
where \tscodeinline{side} is an optional placement --- \tscodeinline{left}, \tscodeinline{right}, \tscodeinline{outer} or \tscodeinline{inner} ---
selecting a margin explicitly. It compiles down to \tscodeinline{\textbackslash{}marginnote\{…\}} from the
\tscodeinline{marginnote} \LaTeX{} package (auto-loaded on first use).
\subsection{Default placement}
No suffix means the document's default side. In a \tscodeinline{oneside} layout that's
the right margin; in a \tscodeinline{twoside} layout, the outer margin (so the note
flips automatically between recto and verso pages).

Most readers appreciate a short quip\tsaside{default note} that sits
beside the paragraph without interrupting the flow. When margin notes are
kept brief and self-contained, they feel like a friendly aside rather than
a distraction.
\subsection{Forced side}
Force the left margin with \tscodeinline{side=left} and the right margin with
\tscodeinline{side=right}. Because
the switch is scoped to a \LaTeX{} group, subsequent unqualified notes still
follow the document's default placement.

Some diagrams read better when labelled on the left\tsaside[side=left]{left-hand
pointer}, while running commentary fits the right-hand
margin\tsaside[side=right]{right-hand commentary} more naturally. Mixing the two in
close succession is fine --- each note is independent.
\subsection{Inline formatting}
Margin notes pass through the full inline-Markdown parser, so they support
the usual \textbf{bold}, \emph{italic}, \tscodeinline{inline code}, and \href{https://example.org}{links}.

Viscoelastic materials can be approximated as linear\tsaside{\textbf{Hooke's}
law applies at small strain where the stress is \emph{proportional} to the
strain} under moderate stress, but non-linear effects dominate above the
yield point.
\subsection{Longer notes}
The package handles multi-sentence notes gracefully. They line-wrap in the
margin and do not disturb the main text's baseline
grid.\tsaside{Longer notes wrap over several lines and keep flowing down
the margin. Keep them concise so the reader can scan them without losing
the thread of the main text.}

Below, a denser example mixing sides and inline styles:
\begin{itemize}
\item Classical mechanics\tsaside[side=left]{see Newton, \emph{Principia}, 1687} predates
the calculus of variations.
\item Modern formulations rely on Lagrangians\tsaside{or Hamiltonians for
energy-based analyses} and variational principles.
\item Quantum mechanics\tsaside[side=right]{introduces non-commuting
observables} follows in the 20th century.
\end{itemize}
\subsection{Inner / Outer}
In a \tscodeinline{twoside} layout --- like this article --- \tscodeinline{\{o\}} (outer) and \tscodeinline{\{i\}} (inner)
are MVP aliases of \tscodeinline{\{r\}} and \tscodeinline{\{l\}} respectively. On single-side documents
they behave identically to their counterparts.
\end{document}
